\begin{afterword}
\summaryhead

The post starts from Drew McDermott's 1976 critique of AI programs that stored facts like ``happiness is a state of mind'' as nodes with suggestive English names. Rename the nodes with random symbols and nothing is left; the names do not grow back. The author argues that a person who believes ``light is waves'' only because a physicist said so holds a belief of the same kind. The test is to ask: ``How would I regenerate this knowledge if it were deleted from my mind?''

Beliefs formed through experience are connected to other knowledge, to perception and to action. Beliefs acquired by being told may not be. The post cites Davidson, by way of Rorty, on beliefs about ``beavers'' so mistaken that they are not about beavers, and Wittgenstein on a wheel that turns nothing. The deeper the imagined deletion, the stricter the test: could the author re-prove the Pythagorean theorem, notice it, or reinvent mathematical proof? A shepherd's apprentice who does not understand the pebble-counting system cannot repair or improve it. The post ends: make yourself the source of every thought worth thinking.

\responsehead

The test the post proposes is a sensible check. Asking whether you could rebuild a piece of knowledge separates understanding from remembering a label, and the shepherd's apprentice shows plainly why the difference matters: without the principle, you cannot adapt the method when circumstances change.

The authorities the post calls on were making other points. McDermott's sentence about getting ``the hearer to notice what it has been told'' concerns programs passing data to each other, a case the paper sets apart from human speech; the post applies it to a person listening to a physicist. The Davidson example concerns beliefs that are mostly false, while the post's case is true beliefs taken on trust. In the article the post cites, Rorty explains the example by the view behind it: a word gets its meaning from shared use among speakers, not from being paired with particular experiences. Wittgenstein's wheel, in its own passage, is the private meaning a word has for a man who ``uses it as we all do.'' The post's argument depends on the suggestion, put as a question, that a belief without enough of the believer's own experience may connect to nothing. On the account its own sources give, a belief is connected through shared use.

The test is not shown in use. The author reports adopting it early, and the Pythagorean example is a pair of guesses (``I think so,'' ``I think I would''). And the fallback offered there, re-proving the theorem after someone else has pointed out the problem and said it has an answer, is much weaker than the original question of rebuilding the knowledge alone. No belief is examined and found to fail the test.

The final advice has no limit. ``Strive to make yourself the source of every thought worth thinking'' comes six weeks after ``Cached Thoughts,'' where the author wrote that ``no one can think fast enough to think their own thoughts.'' The post gives no rule for which thoughts to regenerate and which to accept from others, and that choice is the one a reader faces.

In short: a sensible check on understanding, supported by authorities who were making other points, never shown finding a belief that fails it, and extended to every thought worth thinking with no way to choose among them.
\end{afterword}
